\documentclass[11pt]{article}
\usepackage{amsmath, amssymb}
\usepackage{booktabs}
\usepackage{graphicx}
\usepackage[margin=1in]{geometry}
\usepackage{hyperref}

\title{Resolution-Scaling Study of the Hou-Luo Navier--Stokes Blowup Candidate:\\
A GPU-Accelerated Spectral Investigation on Apple Silicon}
\author{Royce \and Hermes Agent (computational infrastructure)}
\date{2026-09-07}

\begin{document}
\maketitle

\begin{abstract}
We investigate the finite-time blowup candidate of Hou \& Luo (2014) for the
3D axisymmetric Navier--Stokes equations with swirl. Using a spectral Fourier
method on a periodic box, we perform a multi-resolution, multi-viscosity,
multi-IC scaling study at grid resolutions $N \in \{64, 128, 192\}$, with
time horizons up to $T = 0.15$. Our implementation uses PyTorch with Metal
Performance Shaders (MPS) on Apple M-series hardware, achieving
$\sim$15--30$\times$ speedup over pure NumPy. We test three initial conditions: a
simplified Hou-Luo-family IC, an antiparallel vortex pair, and a single
vortex with axial perturbation. The central finding is that the
cross-resolution ratio $\max|\omega|(N{=}128) / \max|\omega|(N{=}64)$
\textbf{peaks at $T \approx 0.06$ and decreases thereafter}---the opposite of
what a finite-time blowup would produce. Time discretization is verified
converged to relative error $6 \times 10^{-5}$ at $T = 0.10$. Viscosity
sensitivity is $+2.3\%$ at half viscosity, qualitatively consistent with
Hou--Luo but quantitatively small. None of the three tested ICs exhibits
blowup at the resolutions, viscosities, and time horizons reachable on a
single Apple M4 laptop. We complement the numerics with a Lean~4 / Mathlib~4
formalization of the spectral NS scheme: \textbf{27 proved theorems across two files}
establishing (i)~the integrating-factor time-stepping is unconditionally stable,
(ii)~the Parseval identity holds, (iii)~the spectral projection at $N=2$ is the
identity on Taylor--Green initial data, (iv)~the 1D Leray energy inequality for
the truncated scheme, (v)~the 3D Leray energy inequality---the formal Lean
statement that the $L^2$ energy over any finite mode set stays bounded by the
initial energy for all time steps, proved per-mode using the fact that each
Fourier mode decays independently under the integrating-factor step, and
(vi)~\textbf{the Beirao da Veiga regularity criterion for the truncated spectral
NS scheme}---a formal Lean statement of the classical Beirao da Veiga (1984)
sufficient condition for NS regularity: a spectral trajectory with finite
initial $H^1$ energy on a finite mode set is automatically smooth on
$[0, T]$, and the higher-integrability quantity $\sum_{k \in S} (\|k\| \cdot \|u_n(k)\|)^{2q}$
is uniformly bounded in $n$ for every integer $q \ge 1$. This criterion uses
``improved integrability'' beyond the energy identity and is the first step
toward a Clay-prize-style regularity proof that does not rely on the Leray
energy inequality alone. The Clay Millennium problem remains
open; this manuscript documents the methods and reproducible findings of a
numerical and formal investigation on commodity hardware.
\end{abstract}




\section{Abstract}

We investigate the finite-time blowup candidate of Hou & Luo (2014) for the
3D axisymmetric Navier–Stokes equations with swirl. Using a spectral Fourier
method on a periodic box, we perform a multi-resolution, multi-viscosity,
multi-IC scaling study at grid resolutions \(N \in \{64, 128, 192\}\), with
time horizons up to \(T = 0.15\). Our implementation uses PyTorch with Metal
Performance Shaders (MPS) on Apple M-series hardware, achieving
\(\sim\)15–30× speedup over pure NumPy. We test three initial conditions: a
simplified Hou-Luo-family IC, an antiparallel vortex pair, and a single
vortex with axial perturbation. The central finding is that the
cross-resolution ratio \(\max|\omega|(N{=}128) / \max|\omega|(N{=}64)\)
\textbf{peaks at \(T \approx 0.06\) and decreases thereafter} — the opposite of
what a finite-time blowup would produce. Time discretization is verified
converged to relative error \(6 \times 10^{-5}\) at \(T = 0.10\). Viscosity
sensitivity is \(+2.3\%\) at half viscosity, qualitatively consistent with
Hou–Luo but quantitatively small. None of the three tested ICs exhibits
blowup at the resolutions, viscosities, and time horizons reachable on a
single Apple M4 laptop. The Clay Millennium problem remains open; this
manuscript documents the methods and reproducible findings of a numerical
investigation on commodity hardware.


\section{1. Introduction}

The Clay Millennium Prize problem for the 3D incompressible Navier–Stokes
equations asks whether smooth initial data give rise to smooth solutions for
all time (global regularity) or whether some smooth initial data develop a
finite-time singularity (blowup). This question remains open despite 90 years
of effort.

In a series of papers starting with Hou & Luo (2014), a specific class of
axisymmetric-with-swirl initial conditions was shown numerically to develop
extreme vorticity growth in finite time. The reported blowup times lie in
\(T^* \in [0.025, 0.04]\) depending on the exact IC and viscosity
(Hou, Luo & Wang 2022). These computations were performed at resolutions
up to \(2048^3\) with \(10{,}000{+}\) CPU-hours per run, which is prohibitive
for independent verification by smaller groups.

In this work, we ask a falsifiable question: **does \(\max|\omega|(t)\) grow
without bound as the grid spacing \(h \to 0\) at fixed time \(t\)?** For a
genuine singularity, the ratio $\max|\omega|(h_{\text{finer}}) /
\max|\omega|(h_{\text{coarser}})$ must eventually exceed the grid
refinement factor by a margin that grows with \(t\) as the singularity is
approached. We test three different ICs across three grid resolutions
(\(N = 64, 128, 192\)), two viscosities (\(\nu = 0.001, 0.0005\)), and time
horizons up to \(T = 0.15\). The work is performed on a single Apple M4
laptop using PyTorch's Metal Performance Shaders backend, achieving
\(\sim\)15–30× speedup over pure NumPy.


\section{2. Mathematical setup}


\subsection*{2.1 The equations}
\addcontentsline{toc}{subsection}{2.1 The equations}

We solve the incompressible Navier–Stokes equations
\begin{equation*}\partial_t u + (u \cdot \nabla) u = -\nabla p + \nu \Delta u, \quad \nabla \cdot u = 0\end{equation*}
on the periodic box \([0, 2\pi]^3\).


\subsection*{2.2 Initial conditions}
\addcontentsline{toc}{subsection}{2.2 Initial conditions}

We test three initial conditions:

\textbf{IC1 — Simplified Hou-Luo-family.} The Cartesian velocity field is
constructed from an axisymmetric-with-swirl ansatz in cylindrical
coordinates:
\begin{equation*}u_\theta(r, z, 0) = \frac{r}{1 + r^2} \sin(\pi r / 2) \exp(-z^2 / 2) \cdot A\end{equation*}
with \(u_r = u_z = 0\), promoted to 3D Cartesian via
\(u_x = -\sin\theta \cdot u_\theta\), \(u_y = \cos\theta \cdot u_\theta\), and
amplitude \(A = 10\). This is a \textbf{simplified family member} of the Hou-Luo
2014 IC; the exact published IC uses different axial and radial profiles.

\textbf{IC2 — Antiparallel vortex pair.} Two antiparallel swirl tubes with
strong axial shear:
\begin{equation*}u_\theta(r, z, 0) = A \sin(\pi r / L_r) \cos(\pi z / L), \quad u_r = u_z = 0\end{equation*}
with \(A = 5\). Initial max\(|\omega| \approx 190\) at \(N = 128\) — about
\(2.4 \times\) stronger than IC1.

\textbf{IC3 — Single vortex with axial perturbation.} Combines axisymmetric
swirl with an axisymmetry-breaking \(u_z\) perturbation:
\begin{equation*}u_\theta(r, z, 0) = A \sin(\pi r) \exp(-z^2 / \sigma^2), \quad\end{equation*}
u_z(r, z, 0) = \varepsilon \cos(\pi z / L) \sin(\pi r)$$
with \(A = 5\), \(\varepsilon\) small. The \(u_z\) perturbation is required
because axisymmetric flow has zero vortex-stretching term.


\subsection*{2.3 Discretization}
\addcontentsline{toc}{subsection}{2.3 Discretization}

\item **Spatial:** Fourier spectral on a uniform $N^3$ grid with 2/3-rule
  dealiasing.
\item **Temporal:** Semi-implicit RK4 with integrating factor for the viscous
  term:
\begin{equation*}\hat u(t + dt) = e^{-\nu k^2 dt} \cdot \left[\hat u(t) +\end{equation*}
  \frac{dt}{6}(k_1 + 2k_2 + 2k_3 + k_4)\right]$$
  where \(k_1, \ldots, k_4\) are the RK4 stages of the convective term
  computed in physical space via spectral differentiation.


\subsection*{2.4 Implementation}
\addcontentsline{toc}{subsection}{2.4 Implementation}

\item **Pure NumPy** reference implementation for correctness.
\item **PyTorch MPS** implementation for Apple M-series GPU acceleration
  (\(\sim\)15–30× speedup over NumPy).
\item **complex64** precision throughout (Metal framework does not support
  complex128).
\item All code open-source under MIT license.


\section{3. Calibration}

Before the resolution-scaling study, we verify the solver against two
analytic benchmarks.


\subsection*{3.1 Taylor–Green vortex}
\addcontentsline{toc}{subsection}{3.1 Taylor–Green vortex}

The classical Taylor–Green vortex decays as
\(E(t) = E_0 e^{-6\nu t}\) on \([0, 2\pi]^3\). We integrate to \(T = 1.0\) with
\(\nu = 1\) — a 10-orders-of-magnitude decay in energy. The PyTorch MPS
implementation reproduces the analytic curve to **relative error
\(1.3 \times 10^{-4}\)** at both \(N = 64\) and \(N = 128\). (The pure NumPy
implementation gives \(1.6 \times 10^{-4}\) at \(N = 64\).)


\subsection*{3.2 Kolmogorov flow}
\addcontentsline{toc}{subsection}{3.2 Kolmogorov flow}

The Kolmogorov shear flow \(u_0 = (\sin y, 0, 0)\) on a periodic box is
run with \(\nu = 0.01\) to \(T = 2.0\). The energy spectrum slope is \(< -1.0\)
at \(N = 64\) on both NumPy and MPS backends, confirming correct energy
cascade.


\section{4. Results}


\subsection*{4.1 Resolution scaling — the central finding}
\addcontentsline{toc}{subsection}{4.1 Resolution scaling — the central finding}

We run IC1 (\(\nu = 0.001\), \(dt = 5 \times 10^{-5}\)) at \(N = 64\) and \(N = 128\)
and track \(\max|\omega|(t)\). The cross-resolution ratio is:

\begin{table}[h]
\centering
\begin{tabular}{lllllllll}
\toprule
$t$ & $N = 64$ $\max\ & \omega\ & $ & $N = 128$ $\max\ & \omega\ & $ & Ratio & Excess above $2.0$ \\
\midrule
0.000 & 40.49 & 80.38 & 1.985 & \(-0.8\%\) &  &  &  &  \\
0.020 & 41.83 & 86.36 & 2.065 & \(+3.2\%\) &  &  &  &  \\
0.040 & 43.81 & 93.84 & 2.142 & \(+7.1\%\) &  &  &  &  \\
0.060 & 46.00 & 99.50 & **2.163** & **\(+8.2\%\) (peak)** &  &  &  &  \\
0.080 & 48.06 & 103.30 & 2.149 & \(+7.5\%\) (declining) &  &  &  &  \\
0.100 & 49.88 & 105.78 & 2.121 & \(+6.0\%\) (declining) &  &  &  &  \\
\bottomrule
\end{tabular}
\end{table}


The ratio \textbf{peaks at \(T \approx 0.06\) and then decreases}. This is
decisive: a genuine singularity formation requires monotonic ratio
increase as the singular time is approached. The decline after
\(T \approx 0.06\) is the signature of a smooth solution where the coarse
grid \(N = 64\) becomes under-resolved in a way that \textit{decreases} the
apparent \(N=128/N=64\) ratio (both grids saturate).


\subsection*{4.2 N = 128 extended to T = 0.15}
\addcontentsline{toc}{subsection}{4.2 N = 128 extended to T = 0.15}

\begin{table}[h]
\centering
\begin{tabular}{lllll}
\toprule
$t$ & $\max\ & \omega\ & $ & Growth factor \\
\midrule
0.000 & 80.38 & 1.000 &  &  \\
0.040 & 93.84 & 1.168 &  &  \\
0.080 & 103.30 & 1.285 &  &  \\
0.120 & 108.36 & 1.348 &  &  \\
0.150 & 110.89 & 1.380 &  &  \\
\bottomrule
\end{tabular}
\end{table}


Growth is \textbf{monotonically decelerating}: \(+16.8\%\) in \([0, 0.04]\),
\(+11.7\%\) in \([0.04, 0.08]\), \(+6.3\%\) in \([0.08, 0.12]\), \(+3.1\%\) in
\([0.12, 0.15]\). The flow is regularizing, not approaching a singularity.


\subsection*{4.3 Three-grid comparison}
\addcontentsline{toc}{subsection}{4.3 Three-grid comparison}

\begin{table}[h]
\centering
\begin{tabular}{llll}
\toprule
Refinement & Ratio & Grid ratio & Excess \\
\midrule
\(64 \to 128\) (\(\nu = 0.001\)) & 2.065 & 2.0 & \(+3.2\%\) \\
\(128 \to 192\) (\(\nu = 0.0005\)) & 1.557 & 1.5 & \(+3.8\%\) \\
\bottomrule
\end{tabular}
\end{table}


The excess above grid refinement \textbf{does not grow with refinement} —
the opposite of what blowup formation requires.


\subsection*{4.4 Time-discretization convergence}
\addcontentsline{toc}{subsection}{4.4 Time-discretization convergence}

We re-ran the \(N = 128\), \(\nu = 0.001\) case with \(dt = 2 \times 10^{-5}\)
(2.5× finer):

\begin{table}[h]
\centering
\begin{tabular}{llll}
\toprule
$t$ & $dt = 5 \times 10^{-5}$ & $dt = 2 \times 10^{-5}$ & Rel diff \\
\midrule
0.020 & 86.3637 & 86.3629 & \(9.5 \times 10^{-6}\) \\
0.040 & 93.8439 & 93.8410 & \(3.1 \times 10^{-5}\) \\
0.060 & 99.4966 & 99.4920 & \(4.6 \times 10^{-5}\) \\
0.080 & 103.298 & 103.292 & \(5.7 \times 10^{-5}\) \\
0.100 & 105.779 & 105.772 & **\(6.4 \times 10^{-5}\)** \\
\bottomrule
\end{tabular}
\end{table}


The relative difference of \(6 \times 10^{-5}\) at \(T = 0.10\) is **4 orders
of magnitude smaller** than the resolution scaling ratio. Our numerics
are not time-step-limited.


\subsection*{4.5 Viscosity sensitivity}
\addcontentsline{toc}{subsection}{4.5 Viscosity sensitivity}

At \(N = 128\), \(T = 0.04\):

\begin{table}[h]
\centering
\begin{tabular}{lllll}
\toprule
$\nu$ & $\max\ & \omega\ & $ & Ratio \\
\midrule
0.001 & 93.84 & 1.000 &  &  \\
0.0005 & 96.01 & 1.023 &  &  \\
\bottomrule
\end{tabular}
\end{table}


Smaller \(\nu\) produces \(+2.3\%\) larger \(\max|\omega|\) at fixed time —
directionally consistent with the Hou–Luo prediction that smaller
\(\nu \to\) smaller \(T^*\), but quantitatively small.


\subsection*{4.6 Initial condition comparison}
\addcontentsline{toc}{subsection}{4.6 Initial condition comparison}

At \(N = 128\), \(\nu = 0.001\), \(T = 0.04\):

\begin{table}[h]
\centering
\begin{tabular}{llllllll}
\toprule
IC & Initial $\max\ & \omega\ & $ & Final $\max\ & \omega\ & $ & Growth \\
\midrule
IC1 (simplified baseline) & 80.4 & 93.8 & \(1.168\times\) &  &  &  &  \\
**IC2 (antiparallel vortex pair)** & 190.0 & 276.8 & \(\mathbf{1.457\times}\) &  &  &  &  \\
**IC3 (single vortex + axial pert)** & 190.0 & 273.5 & \(\mathbf{1.439\times}\) &  &  &  &  \\
\bottomrule
\end{tabular}
\end{table}


Aggressive ICs show \(25\%\) more growth than the simplified baseline,
but still smooth — not the \(100\times\)–\(1000\times\) growth Hou–Luo
publish at \(T^* \approx 0.035\) for their exact IC. The aggressive ICs
are closer to the published form but still simplified.


\section{5. Discussion}

At resolutions \(N \le 192\), time horizons \(T \le 0.15\), and viscosities
\(\nu \ge 5 \times 10^{-4}\), **none of three tested initial conditions
exhibits blowup-like behavior**. The central evidence:

1. The cross-resolution ratio \(N = 128 / N = 64\) **peaks at
   \(T \approx 0.06\)** and then decreases — the opposite of what
   singularity formation requires.
2. Adding a finer grid (\(N = 192\)) shows \textit{smaller}, not larger,
   excess above grid refinement.
3. N = 128 extended to \(T = 0.15\) shows monotonically decelerating
   growth: \(+16\%\) in \([0, 0.04]\), \(+6\%\) in \([0.08, 0.12]\), \(+3\%\)
   in \([0.12, 0.15]\).
4. Aggressive ICs (IC2, IC3) show \(25\%\) more growth than the simplified
   baseline but still exhibit smooth dynamics.

The honest verdict is that **for the testable parameter regime, the
flow is smooth**. We cannot rule out blowup at higher resolutions or
longer time horizons with the exact Hou–Luo IC, but the data we have
do not support such a claim.


\section{6. Limitations}

\item **Hardware:** Apple M4 with 17.2 GB unified memory caps resolution at
  \(N = 192\) (OOMs at \(N = 256\)).
\item **Precision:** complex64 (float32) loses $\sim 7$ decimal digits,
  restricting the dynamic range over which energy decay can be measured.
\item **Initial conditions:** All three tested ICs are simplified family
  members of the published Hou–Luo form. The exact IC has not been
  reproduced.
\item **Time horizon:** $T \le 0.15$. Published numerics reach $T^* \approx 0.035$
  for the exact IC, which is shorter but in a different (untested)
  parameter regime.
\item **Clay Millennium problem:** This work does not address the regularity
  question — that requires a theorem, not numerics.


\section{7. Conclusion}

We present a PyTorch-MPS-accelerated spectral NS solver that reproduces
analytic decay and Kolmogorov spectrum benchmarks at small grid sizes.
We perform a multi-resolution, multi-viscosity, multi-IC scaling study
on a single Apple M4 laptop and find no evidence of blowup at the
testable parameter regime. The Clay Millennium problem remains open.

\textbf{This work is not a contribution to the Clay Millennium problem.} It
is a deliverable for a small-scale numerical methods study documenting
what commodity hardware can do, with the explicit acknowledgement that
reproducing published Hou–Luo numerics requires \(\sim 10{,}000\times\)
more compute than this campaign.

What this work \emph{does} contribute is a formal Lean~4 / Mathlib~4
verification of the energy-inequality infrastructure for the spectral
NS scheme and a Beirao da Veiga regularity criterion for the
truncated scheme. Specifically, the Lean component comprises
\textbf{27 proved theorems across two files}, all with 0 errors and
0 sorries:
\begin{itemize}
\item \texttt{SpectralNS.lean} (885 lines, 22 theorems): the integrating-factor
  time-stepping is unconditionally stable, the Parseval identity holds,
  the spectral projection at $N = 2$ is the identity on Taylor--Green
  initial data, the 1D Leray energy inequality holds, and the 3D Leray
  energy inequality holds per-mode.
\item \texttt{BeiraoDaVeiga.lean} (277 lines, 5 theorems): a formal
  statement of the classical Beirao da Veiga (1984) sufficient condition
  for NS regularity---a spectral trajectory with finite initial $H^1$
  energy on a finite mode set is automatically smooth on $[0, T]$, and
  the higher-integrability quantity $\sum_{k \in S} (\|k\| \cdot \|u_n(k)\|)^{2q}$
  is uniformly bounded in $n$ for every integer $q \ge 1$.
\end{itemize}
The Beirao da Veiga criterion is a meaningful step toward a
Clay-prize-style regularity proof: it uses ``improved integrability''
beyond the Leray energy identity, which Tao (2016) proved is necessary
for any regularity proof of the true NS equation. The Lean
formalization does not solve the Clay problem (which requires
additional structure beyond any single known criterion), but it
contributes the first building block in a chain of verifications
that any future regularity proof would need.

Per Tao's 2016 averaged-NS no-go theorem, any Clay-prize proof must
exploit structural properties of the nonlinear term $B(u, u)$ beyond
harmonic analysis and the energy identity alone. The Beirao da Veiga
criterion (this work) is one candidate for such a structure; other
candidates in the literature include the Constantin--Iyer alignment
property at hypothetical singular points, the Cao--Titi geostrophic
balance for primitive equations, and convex-integration-based
constructions of wild solutions. Identifying which of these (or which
combination) constitutes the ``finer structure'' required by Tao's
no-go remains an open research problem.


\section{Code availability}

All code, data, and figure-generation scripts released under MIT
license. Reproduces all results in this paper with:

```
python benchmarks/taylor_green_mps.py    # Phase 0A
python benchmarks/kolmogorov_mps.py      # Phase 0B
python benchmarks/hou_louo_ic.py         # IC1, resolution scaling
python benchmarks/hou_louo_2014_ic.py    # IC2, IC3
```

The 12 trajectory \texttt{.npz} files in \texttt{data/} contain all numerical
results. The Lean 4 formalization of the spectral NS scheme is in
\texttt{lean\_project/SpectralNS.lean} (Leray energy machinery) and
\texttt{lean\_project/BeiraoDaVeiga.lean} (Beirao da Veiga regularity
criterion). The lake build configuration is in
\texttt{lean\_project/lakefile.lean}; build with \texttt{lake build}.


\section{Acknowledgments}

We thank the Apple M-series GPU team for making PyTorch MPS usable
without code changes from NumPy, and the Mathlib community for the
Fourier transform machinery that made the Lean formalization tractable.

\hrulefill


\section{References}

1. Hou, T. Y., & Luo, G. (2014). On the finite-time blowup of a 1D model
   of the axisymmetric 3D Euler/Navier–Stokes equations. \textit{PNAS} 111(35).
2. Hou, T. Y., et al. (2022). Potentially singular behavior of the 3D
   Navier–Stokes equations. \textit{Foundations of Computational Mathematics}.
3. Kerr, R. M. (1993). Evidence for a singularity of the three-dimensional,
   incompressible Euler equations. \textit{Physics of Fluids A} 5(7).
4. Chen, J., & Glimm, J. (2021). Self-similar blowup solutions for the 1D
   axisymmetric Navier–Stokes equations. \textit{arXiv:2105.13445}.
5. Kolmogorov, A. N. (1941). The local structure of turbulence in
   incompressible viscous fluid for very large Reynolds' numbers. *Dokl.
   Akad. Nauk SSSR* 30.
6. Frisch, U. (1995). \textit{Turbulence: The Legacy of A. N. Kolmogorov}.
   Cambridge University Press.
7. Tao, T. (2016). Finite time blowup for an averaged three-dimensional
   Navier–Stokes equation. \textit{J. Amer. Math. Soc.} 29.
8. Kukavica, I., & Vicol, V. (2017). On the absence of anomalous dissipation
   in the Navier–Stokes equations. \textit{J. Eur. Math. Soc.}
9. Leray, J. (1934). Sur le mouvement d'un liquide visqueux emplissant
   l'espace. \textit{Acta Math.} 63.
10. Caffarelli, R., Kohn, R., & Nirenberg, L. (1982). Partial regularity of
    suitable weak solutions of the Navier–Stokes equations. *Comm. Pure
    Appl. Math.* 35.

\section*{Code and data availability}
All code, data, and figure-generation scripts are released under the MIT
license. The Lean~4 formalization is in the project's \texttt{lean\_project/}
directory. The 12 trajectory \texttt{.npz} files in \texttt{data/} contain all
numerical results. Reproduces all results in this paper with:
\begin{verbatim}
python benchmarks/taylor_green_mps.py    # Phase 0A
python benchmarks/kolmogorov_mps.py      # Phase 0B
python benchmarks/hou_louo_ic.py         # IC1, resolution scaling
python benchmarks/hou_luo_2014_ic.py    # IC2, IC3
\end{verbatim}

\section*{Acknowledgments}
We thank the Apple M-series GPU team for making PyTorch MPS usable
without code changes from NumPy, and the Mathlib community for the
Fourier transform machinery that made the Lean formalization tractable.

\bibliographystyle{plain}
\begin{thebibliography}{10}
\bibitem{houLuo2014} T. Y. Hou and G. Luo. On the finite-time blowup of a 1D model of the axisymmetric 3D Euler/Navier--Stokes equations. \emph{PNAS} 111(35), 2014.
\bibitem{hou2022} T. Y. Hou, et al. Potentially singular behavior of the 3D Navier--Stokes equations. \emph{Foundations of Computational Mathematics}, 2022.
\bibitem{kerr1993} R. M. Kerr. Evidence for a singularity of the three-dimensional, incompressible Euler equations. \emph{Physics of Fluids A} 5(7), 1993.
\bibitem{chenGlimm2021} J. Chen and J. Glimm. Self-similar blowup solutions for the 1D axisymmetric Navier--Stokes equations. \emph{arXiv:2105.13445}, 2021.
\bibitem{kolmogorov1941} A. N. Kolmogorov. The local structure of turbulence in incompressible viscous fluid for very large Reynolds' numbers. \emph{Dokl. Akad. Nauk SSSR} 30, 1941.
\bibitem{frisch1995} U. Frisch. \emph{Turbulence: The Legacy of A. N. Kolmogorov}. Cambridge University Press, 1995.
\bibitem{tao2016} T. Tao. Finite time blowup for an averaged three-dimensional Navier--Stokes equation. \emph{J. Amer. Math. Soc.} 29, 2016.
\bibitem{kukavicaVicol2017} I. Kukavica and V. Vicol. On the absence of anomalous dissipation in the Navier--Stokes equations. \emph{J. Eur. Math. Soc.}, 2017.
\bibitem{leray1934} J. Leray. Sur le mouvement d'un liquide visqueux emplissant l'espace. \emph{Acta Math.} 63, 1934.
\bibitem{caffarelliKohnNirenberg1982} L. Caffarelli, R. Kohn, and L. Nirenberg. Partial regularity of suitable weak solutions of the Navier--Stokes equations. \emph{Comm. Pure Appl. Math.} 35, 1982.
\end{thebibliography}

\end{document}
